%!TEX root = ../main.tex

\chapter{Metodología y procedimientos}\label{cap:metodologia}

\section{Metodología de la investigación}

\subsection{Área temática}

La presente investigación se desarrolla en el campo de la Ingeniería en Informática y está orientada al diseño de una solución tecnológica para asistir a cuidadores de personas con discapacidad.

El proyecto combina conocimientos relacionados con el desarrollo de software, la gestión de información y los sistemas de recomendación, con el objetivo de brindar apoyo en tareas que forman parte del cuidado diario. Entre ellas se encuentran la organización de turnos médicos, el seguimiento de la medicación, la coordinación entre quienes cuidan, la realización de trámites y el acceso a recursos útiles para cada situación particular.

Asimismo, la propuesta se relaciona con el ámbito de la discapacidad, ya que busca dar respuesta a una problemática real mediante el uso de herramientas tecnológicas que faciliten el acceso a la información y contribuyan a una mejor toma de decisiones por parte de los cuidadores.

\subsection{Planteamiento del problema}

El cuidado de una persona con discapacidad implica una gran cantidad de responsabilidades que, en la mayoría de los casos, recaen sobre familiares que no cuentan con formación específica en cuestiones médicas, administrativas o de gestión. Además de acompañar al paciente en sus actividades diarias, los cuidadores deben organizar consultas médicas, administrar tratamientos, gestionar documentación, realizar trámites y buscar recursos adecuados para cada situación.

Actualmente, gran parte de la información necesaria para desarrollar estas tareas se encuentra distribuida entre diferentes organismos, centros de salud, páginas web, grupos informales de comunicación y documentación física. Esta dispersión obliga a los cuidadores a invertir tiempo y esfuerzo en la búsqueda de información, muchas veces sin la certeza de que esta sea correcta o se encuentre actualizada.

Como consecuencia, es frecuente que se produzcan retrasos en la realización de trámites, pérdida de turnos, dificultades para acceder a determinados servicios y una mayor carga emocional para quienes ejercen el rol de cuidador. Cuando el cuidado se reparte entre varias personas, como un familiar y un cuidador contratado, o entre los empleados de una institución, se agrega la dificultad de saber qué se hizo y quién lo hizo. A esto se suma la ausencia de herramientas tecnológicas que permitan integrar en un único entorno las distintas necesidades asociadas al acompañamiento de una persona con discapacidad.

\subsection{Delimitación de la investigación}

Esta investigación se enfoca en estudiar cómo una herramienta informática puede contribuir a mejorar la experiencia de las personas que asumen tareas de cuidado. El trabajo busca comprender las principales dificultades que enfrentan los cuidadores en su vida cotidiana y desarrollar una propuesta tecnológica que les brinde acompañamiento y apoyo en distintas situaciones relacionadas con esta responsabilidad.

El proyecto se centra en la creación de una plataforma que reúna recursos, información y herramientas de asistencia orientadas a facilitar la organización de actividades, la coordinación del grupo de cuidado y el acceso a orientación práctica. Además, se analiza el uso de técnicas de sistemas de recomendación como complemento para ofrecer sugerencias adaptadas a las necesidades de cada persona.

La propuesta está dirigida principalmente a cuidadores no profesionales, quienes muchas veces desempeñan esta labor sin contar con capacitación específica ni espacios de contención adecuados. También contempla a quienes cuidan de forma remunerada, como cuidadores contratados por una familia o empleados de una institución, y a las familias que necesitan seguir ese cuidado a distancia. Por este motivo, la investigación pone el foco en sus necesidades, experiencias y desafíos diarios.

La propuesta no busca reemplazar el trabajo de médicos, psicólogos u otros profesionales, sino brindar una herramienta de apoyo y orientación para las personas cuidadoras. Por este motivo, no se incluyen funcionalidades relacionadas con diagnósticos, indicaciones de tratamiento ni decisiones clínicas.

Como resultado, se espera obtener una herramienta que permita validar la propuesta planteada y analizar de qué manera la tecnología puede contribuir a mejorar la experiencia y la calidad de vida de las personas que ejercen tareas de cuidado.

\subsection{Marco teórico}

El marco teórico de esta investigación está compuesto por los conceptos, enfoques y antecedentes que permiten comprender tanto la problemática de las personas cuidadoras como las posibilidades que ofrecen las tecnologías actuales para brindarles apoyo. Para ello, se estudian diferentes perspectivas relacionadas con el cuidado informal, el impacto que esta actividad tiene en la vida de quienes la desempeñan y las alternativas existentes para acompañar y facilitar esta tarea.

En primer lugar, se analiza el rol de las personas cuidadoras dentro del entorno familiar y social, considerando las responsabilidades que asumen en el acompañamiento de personas con discapacidad. Se busca comprender sus características, las dificultades que enfrentan en su día a día y las necesidades que suelen surgir como consecuencia de las exigencias propias del cuidado.

Asimismo, se profundiza en conceptos relacionados con el bienestar y la calidad de vida de las personas cuidadoras. Se estudian aspectos vinculados a la carga física y emocional que puede generar esta labor, así como las consecuencias que pueden derivarse de la falta de apoyo, información o espacios de contención. También se revisan investigaciones y experiencias previas que permiten comprender la importancia de promover herramientas que contribuyan al acompañamiento y la orientación de este grupo.

Por otra parte, se examina el papel que desempeñan las tecnologías digitales en la creación de soluciones orientadas al apoyo de las personas. En este sentido, se analizan diferentes tipos de plataformas, aplicaciones y servicios digitales que facilitan el acceso a información, recursos y asistencia, identificando aquellas características que resulten relevantes para el desarrollo de la propuesta.

Desde el enfoque tecnológico, se estudian los fundamentos de los sistemas de recomendación y de la personalización, y se revisan también los sistemas conversacionales con inteligencia artificial, como alternativa para brindar orientación. Se analizan las posibilidades y los riesgos de cada enfoque, siempre entendiendo la tecnología como complemento del acompañamiento humano y profesional.

Finalmente, se abordan conceptos relacionados con la experiencia de usuario, la accesibilidad y el diseño centrado en las personas. Se analizan criterios de usabilidad, claridad de la información y facilidad de interacción, con el propósito de desarrollar una solución que pueda ser utilizada de manera simple e intuitiva por usuarios con distintos niveles de conocimiento tecnológico. Los conocimientos obtenidos a partir de este análisis sirven como fundamento para definir las características y funcionalidades de la solución propuesta.

\subsection{Diseño concreto de investigación}

El diseño de la investigación está orientado a comprender las necesidades, dificultades y experiencias de las personas que desempeñan tareas de cuidado, con el propósito de identificar oportunidades en las que una solución tecnológica pueda aportar apoyo y acompañamiento. Para ello, se realiza una revisión de antecedentes, estudios e iniciativas relacionadas con el cuidado informal, así como de herramientas digitales existentes que busquen asistir a este grupo de usuarios.

Con el fin de aproximarse a la realidad que se pretende abordar, se plantean encuestas y entrevistas dirigidas a personas cuidadoras. Estas técnicas permiten relevar información sobre las principales problemáticas que enfrentan, los recursos que utilizan actualmente, las dificultades que encuentran para acceder a información o apoyo, y las expectativas que podrían tener respecto de una herramienta tecnológica orientada a sus necesidades.

La estrategia de investigación combina enfoques cuantitativos y cualitativos, en un diseño mixto \citep{hernandezsampieri2014}. Por un lado, las encuestas permiten obtener información general sobre las características y necesidades de la población objetivo. Por otro, las entrevistas brindan una comprensión más profunda de las experiencias, percepciones y desafíos que forman parte de la vida cotidiana de las personas cuidadoras. A esto se suman las pruebas técnicas del prototipo y la evaluación del sistema de recomendación con datos de demostración.

\subsection{Indicadores}

Para evaluar el avance y la efectividad de la propuesta, se definen indicadores orientados a medir tanto la utilidad percibida de la solución como su capacidad de respuesta frente a las necesidades identificadas en las personas cuidadoras.

Se consideran indicadores vinculados al nivel de comprensión y facilidad de uso de la plataforma, con el objetivo de analizar si las personas cuidadoras pueden interactuar con la solución de manera intuitiva y sin dificultades relevantes. También se tiene en cuenta la utilidad percibida de la información y el apoyo brindado, evaluando si los usuarios encuentran valor en las funcionalidades propuestas y si estas responden a situaciones reales de su día a día.

Otro aspecto para medir es el nivel de satisfacción general con la herramienta, considerando la percepción del usuario respecto a la claridad de la información, la relevancia de las recomendaciones y la sensación de acompañamiento que pueda generar el sistema. Adicionalmente, se evalúa la intención de uso futuro, la predisposición de los usuarios a volver a utilizar la plataforma o recomendarla a otras personas en situaciones similares.

Para el sistema de recomendación se definen además indicadores propios. La proporción de recomendaciones que incluyen una explicación, la cobertura del catálogo y, una vez que haya usuarios reales, la precisión, entendida como la proporción de recursos marcados como útiles entre los primeros recomendados.

\subsection{Técnicas de recolección de datos}

Las técnicas de recolección de datos se implementan a partir del diseño de investigación definido previamente, con el objetivo de obtener información tanto del contexto teórico como de la experiencia real de las personas cuidadoras.

Se realiza una revisión de antecedentes y fuentes bibliográficas relacionadas con el cuidado informal y las tecnologías de apoyo, complementada con el análisis de soluciones digitales existentes orientadas a problemáticas similares.

Por otra parte, se plantean encuestas dirigidas a personas cuidadoras, así como entrevistas en profundidad a un grupo reducido de participantes, lo que permite contrastar información general con experiencias más detalladas.

Finalmente, se incluyen instancias de validación mediante pruebas de uso del prototipo y el registro de las valoraciones que los cuidadores hacen de cada recurso, con el fin de observar su funcionamiento en situaciones reales y obtener retroalimentación directa de los usuarios.

\subsection{Instrumentos de recolección de datos}

Para la recolección de datos se utilizan distintos instrumentos que permiten aplicar las técnicas definidas previamente de manera organizada y consistente.

Para la realización de encuestas se utiliza un formulario digital de doce preguntas, la mayoría cerradas y dos de respuesta múltiple, que relevan el vínculo con la persona cuidada, el tiempo y las horas dedicadas al cuidado, el tipo de discapacidad, la dificultad para organizar la información, los problemas más frecuentes, las herramientas que se usan hoy, el nivel de estrés y la disposición a usar una plataforma como la propuesta. Una última pregunta abierta indaga qué es lo que más ayudaría en la tarea de cuidar.

\subsection{Procesamiento de los datos}

Una vez recopilada la información, los datos obtenidos mediante encuestas, entrevistas y pruebas con usuarios son organizados y clasificados para facilitar su análisis. Los resultados de las encuestas permiten identificar tendencias y necesidades comunes, mientras que las entrevistas aportan una visión más profunda sobre las experiencias y desafíos de las personas cuidadoras.

Luego, se analiza la información obtenida con el objetivo de evaluar si la solución propuesta responde a las necesidades detectadas durante la investigación. Este análisis permite identificar fortalezas, aspectos a mejorar y oportunidades para futuras ampliaciones de la herramienta.

\subsection{Análisis de los datos}

Una vez analizada la información recopilada, se busca comprender de qué manera las necesidades y experiencias relevadas se relacionan con la solución propuesta. Para ello, se toman como referencia los indicadores definidos previamente y los objetivos establecidos para la investigación.

Este análisis permite determinar si la herramienta desarrollada resulta útil para las personas cuidadoras y si responde a las problemáticas identificadas durante el relevamiento. Además, ayuda a detectar posibles mejoras y aspectos a considerar en futuras etapas de desarrollo.

\subsection{Síntesis y conclusiones}

A partir de los resultados obtenidos durante la investigación, se elaboran las conclusiones del proyecto. Estas permiten determinar en qué medida la solución propuesta responde a las necesidades identificadas en las personas cuidadoras y si los objetivos planteados fueron alcanzados.

Las conclusiones sirven para reconocer los principales aportes de la propuesta, identificar sus limitaciones y señalar posibles líneas de mejora o ampliación para futuros desarrollos. De esta manera, se busca obtener una visión integral sobre el potencial de la herramienta como medio de apoyo y acompañamiento para las personas que desempeñan tareas de cuidado.

\section{Metodología de dirección del proyecto}

\subsection{Marco de trabajo}

El proyecto se dirigió con un enfoque iterativo e incremental, alineado con los principios de la guía del PMBOK \citep{pmi2021}. Cada entrega de la cátedra funcionó como una iteración con un entregable concreto, un capítulo del documento y, a partir del capítulo 4, un incremento del sistema que se probaba antes de avanzar. Se eligió este enfoque, en lugar de un plan en cascada, porque el alcance se fue ajustando con las devoluciones de la cátedra, por ejemplo, al reemplazar la comunidad de cuidadores por la sección de artículos y al rediseñar los grupos y los planes. Al tratarse de una sola persona, no se adoptaron las ceremonias de un marco ágil como Scrum, pero sí sus prácticas centrales, trabajo en incrementos cortos, una lista priorizada de pendientes y revisión al cierre de cada iteración.

\subsection{Alcance}

El proyecto abarca el análisis, diseño y desarrollo de una plataforma digital orientada a brindar apoyo y orientación a personas cuidadoras, a las familias que siguen ese cuidado y a las instituciones que lo prestan. Incluye el relevamiento de necesidades, el desarrollo de un prototipo funcional y su validación.

No forman parte del alcance las funcionalidades relacionadas con diagnósticos médicos, indicaciones de tratamiento ni historias clínicas. Para las instituciones se incluye una vista de supervisión del cuidado, pero no su gestión administrativa, como la facturación o la organización del personal. Tampoco se incluyen foros abiertos entre usuarios, por el riesgo de que circulen consejos sobre medicación o tratamientos.

\subsection{Tiempos}

El proyecto se desarrolla en distintas etapas que incluyen la investigación de la problemática, el relevamiento de información, el análisis de los datos obtenidos, el diseño y desarrollo de la solución propuesta, su validación y la elaboración de la documentación final. La planificación de estas actividades se representa en el diagrama de Gantt de la figura~\ref{fig:gantt}. Los hitos de las entregas 5 a 8 y del cierre corresponden al cronograma de la cátedra, y el inicio de mayo es aproximado.

\begin{figure}[tbp]
\centering
\footnotesize
\begin{ganttchart}[
    time slot format=isodate,
    time slot unit=day,
    x unit=0.047cm,
    y unit chart=0.48cm,
    y unit title=0.5cm,
    hgrid, vgrid={*6{draw=none},dotted},
    title label font=\scriptsize,
    bar label font=\scriptsize,
    milestone label font=\scriptsize\itshape,
    bar/.append style={fill=mainColor!60},
    milestone/.append style={fill=mainColor},
    bar height=0.55,
    milestone height=0.55,
  ]{2026-05-04}{2026-11-20}
  \gantttitlecalendar{month=shortname} \\
  \ganttbar{Capítulos 1 a 4}{2026-05-04}{2026-09-03} \\
  \ganttbar{Desarrollo del sistema}{2026-08-17}{2026-11-03} \\
  \ganttbar{Capítulo 5}{2026-09-04}{2026-09-15} \\
  \ganttmilestone{Entrega 5}{2026-09-15} \\
  \ganttbar{Capítulo 6}{2026-09-16}{2026-09-29} \\
  \ganttmilestone{Entrega 6}{2026-09-29} \\
  \ganttbar{Relevamiento y capítulo 7}{2026-09-30}{2026-10-13} \\
  \ganttmilestone{Entrega 7}{2026-10-13} \\
  \ganttbar{Capítulo 8}{2026-10-14}{2026-10-20} \\
  \ganttmilestone{Entrega 8}{2026-10-20} \\
  \ganttbar{Integración del documento}{2026-10-21}{2026-10-27} \\
  \ganttbar{Preparación del despliegue y congelamiento}{2026-10-21}{2026-11-03} \\
  \ganttmilestone{Sistema congelado}{2026-11-03} \\
  \ganttbar{Presentación y ensayo}{2026-11-04}{2026-11-10} \\
  \ganttmilestone{Entrega final}{2026-11-17}
\end{ganttchart}
\caption{Cronograma del proyecto}\label{fig:gantt}
\end{figure}

\subsection{Costos}

El desarrollo del proyecto no contempla una inversión económica significativa, ya que se utilizan herramientas de software gratuitas o de acceso académico. Los principales recursos necesarios están asociados al uso de plataformas digitales para la recolección de datos, el desarrollo del prototipo y la documentación del trabajo.

En caso de realizar entrevistas o validaciones con usuarios, los costos son mínimos y están relacionados principalmente con conectividad o tiempo de los participantes, sin requerir infraestructura adicional ni equipamiento especializado.

\subsection{Riesgos}

Durante el desarrollo del proyecto pueden presentarse algunos riesgos que afecten el cumplimiento de los plazos o la calidad de los resultados. Entre ellos se considera la posible dificultad para acceder a personas cuidadoras que participen en encuestas o entrevistas, lo que podría limitar la cantidad de información disponible.

También se contempla la posibilidad de demoras en la recolección o análisis de datos, así como dificultades técnicas durante el desarrollo del prototipo o la validación con usuarios.

Como estrategia de mitigación, hay una planificación anticipada de las actividades de recolección de datos, se utilizan canales digitales para ampliar el alcance de la muestra y se realizan pruebas progresivas que permiten detectar errores en etapas tempranas del desarrollo.

\subsection{Calidad}

La calidad del proyecto se garantiza mediante la validación progresiva de cada una de sus etapas, asegurando coherencia entre los objetivos planteados, la información recolectada y la solución desarrollada.

Se contemplan revisiones periódicas del avance del trabajo, pruebas automatizadas del sistema, la verificación de las fuentes de información contra páginas oficiales y la validación de los resultados obtenidos a partir de encuestas, entrevistas y pruebas con usuarios. Además, se tiene en cuenta la retroalimentación de los participantes para identificar posibles mejoras en la propuesta.

\subsection{Recursos}

El proyecto requiere recursos humanos, tecnológicos y de información. Dentro de los recursos humanos se incluye la participación de personas cuidadoras para encuestas, entrevistas y validaciones del prototipo.

En cuanto a los recursos tecnológicos, se utilizan computadoras personales, herramientas de desarrollo de software, plataformas digitales para recolección de datos y servicios en línea para documentación y organización del trabajo.

Como recursos de información se consideran bibliografías, artículos académicos, normativas y fuentes digitales oficiales relacionadas con el cuidado informal, la discapacidad, la salud digital y el uso de tecnologías de apoyo.

\subsection{Satisfacción}

La satisfacción de los usuarios se evalúa a partir de la percepción que tengan las personas cuidadoras sobre la utilidad, claridad y facilidad de uso de la herramienta desarrollada.

Se considera el grado en que la solución logra brindar apoyo en situaciones cotidianas de cuidado, así como la disposición de los usuarios a volver a utilizarla o recomendarla. También se tiene en cuenta la sensación de acompañamiento que la plataforma pueda generar y, en el caso de las familias que contratan a alguien y de las instituciones, la tranquilidad de saber cómo se está cuidando a cada persona.

\section{Declaración sobre el uso de inteligencia artificial}

Durante el proyecto utilicé Claude, el asistente de inteligencia artificial de Anthropic, a través de Claude Code, como apoyo en el desarrollo del sistema, la estructura del documento, la evaluación económica y algunas búsquedas bibliográficas. Las decisiones de alcance, diseño y funcionamiento del sistema, los criterios de evaluación y las conclusiones fueron definidos por mí. El código y los resultados se revisaron y verificaron antes de incorporarse al trabajo, y las cifras y referencias se contrastaron con sus fuentes originales.

